% GENERATED FILE. Edit canonical lessons, then run npm run generate:books.
% canonical-source-hash: fnv1a64:a690dcb8e0d59688
% canonical-lessons: PA-C345-sharat, PA-C345-faisla, PA-C345-rae, PA-C345-mukauna, PA-C345-tajviz

\chapter{Making Up Our Minds}
\label{ch:pa-a2-making-up-our-minds}
\hlchaptermodality{345}

\begin{chapteropening}
\textbf{By the end of this chapter:} \emph{I can set a condition and say what we decided.}

Complete the last lesson of chapter 345: I can set a condition and say what we decided.
\end{chapteropening}

\section[\texorpdfstring{\pa{ਸ਼ਰਤ} (sharat)}{ਸ਼ਰਤ (sharat)}]{\pa{ਸ਼ਰਤ} (sharat) — a condition}

\label{lesson:PA-C345-sharat}



\begin{quote}
Before the new one: say the Punjabi for stubbornness, then the Punjabi for an objection.
\end{quote}

\subsection*{You'll want to know: \pa{ਸ਼ਰਤ}}
\textbf{\pa{ਸ਼ਰਤ}} — \emph{sharat} — \textquotedblleft{}a condition\textquotedblright{}.

\begin{quote}
\pa{ਮੇਰੀ} \pa{ਇੱਕ} \pa{ਸ਼ਰਤ} \pa{ਹੈ} \pa{ਅਤੇ} \pa{ਉਹ} \pa{ਔਖੀ} \pa{ਨਹੀਂ} \pa{ਹੈ।}
\end{quote}

\begin{quote}
\emph{merī ikk sharat hai ate oh aukhī nahīṁ hai.}
\end{quote}

\begin{quote}
I have one condition, and it is not hard.
\end{quote}

A feminine noun.

\begin{cousinweb}
\textbf{\pa{ਸ਼ਰਤ}} (\emph{sharat}) is a loan, usually traced through Persian to Arabic \emph{sharṭ}, \textquotedblleft{}condition\textquotedblright{}.
\end{cousinweb}

\begin{grammarlens}[title={What you've built}]
One more word for asking and answering.
\end{grammarlens}

\subsection*{Your turn}
\begin{itemize}
\raggedright
  \item \emph{Say it:} \emph{sharat}
  \item \emph{Say it:} \emph{sharat}, once more
  \item \emph{Say it:} \emph{merī ikk sharat hai ate oh aukhī nahīṁ hai.} — I have one condition, and it is not hard.
  \item \emph{Recall:} say the Punjabi for stubbornness, then the Punjabi for an objection, then say \emph{sharat} again
\end{itemize}

\begin{tcolorbox}[breakable,colback=teal!4,colframe=teal!35!black,title={Before you move on}]
What does \pa{ਸ਼ਰਤ} mean? (\textquotedblleft{}A condition\textquotedblright{}.) Say it once more.
\end{tcolorbox}
\section[\texorpdfstring{\pa{ਫ਼ੈਸਲਾ} (faislā)}{ਫ਼ੈਸਲਾ (faislā)}]{\pa{ਫ਼ੈਸਲਾ} (faislā) — a decision}

\label{lesson:PA-C345-faisla}



\begin{quote}
Before the new one: say the Punjabi for an objection, then the Punjabi for a condition.
\end{quote}

\subsection*{You'll want to know: \pa{ਫ਼ੈਸਲਾ}}
\textbf{\pa{ਫ਼ੈਸਲਾ}} — \emph{faislā} — \textquotedblleft{}a decision\textquotedblright{}.

\begin{quote}
\pa{ਕੀ} \pa{ਤੁਹਾਡਾ} \pa{ਫ਼ੈਸਲਾ} \pa{ਪੱਕਾ} \pa{ਹੈ}?
\end{quote}

\begin{quote}
\emph{kī tuhāḍā faislā pakkā hai?}
\end{quote}

\begin{quote}
Is your decision final?
\end{quote}

A masculine noun.

\begin{cousinweb}
\textbf{\pa{ਫ਼ੈਸਲਾ}} (\emph{faislā}) is a loan, usually traced through Persian to Arabic \emph{faiṣala}, \textquotedblleft{}decision\textquotedblright{}.
\end{cousinweb}

\begin{grammarlens}[title={What you've built}]
One more word for asking and answering.
\end{grammarlens}

\subsection*{Your turn}
\begin{itemize}
\raggedright
  \item \emph{Say it:} \emph{faislā}
  \item \emph{Say it:} \emph{faislā}, once more
  \item \emph{Say it:} \emph{kī tuhāḍā faislā pakkā hai?} — Is your decision final?
  \item \emph{Recall:} say the Punjabi for an objection, then the Punjabi for a condition, then say \emph{faislā} again
  \item \emph{Write it:} copy \textbf{\pa{ਫ਼ੈਸਲਾ}} once, checking the whole word
\end{itemize}

\begin{tcolorbox}[breakable,colback=teal!4,colframe=teal!35!black,title={Before you move on}]
What does \pa{ਫ਼ੈਸਲਾ} mean? (\textquotedblleft{}A decision\textquotedblright{}.) Say it once more.
\end{tcolorbox}
\section[\texorpdfstring{\pa{ਰਾਏ} (rāe)}{ਰਾਏ (rāe)}]{\pa{ਰਾਏ} (rāe) — an opinion; a view}

\label{lesson:PA-C345-rae}



\begin{quote}
Before the new one: say the Punjabi for a condition, then the Punjabi for a decision.
\end{quote}

\subsection*{You'll want to know: \pa{ਰਾਏ}}
\textbf{\pa{ਰਾਏ}} — \emph{rāe} — \textquotedblleft{}an opinion; a view\textquotedblright{}.

\begin{quote}
\pa{ਕੀ} \pa{ਤੁਸੀਂ} \pa{ਮੇਰੀ} \pa{ਰਾਏ} \pa{ਮੰਨਦੇ} \pa{ਹੋ}?
\end{quote}

\begin{quote}
\emph{kī tusī̃ merī rāe mannde ho?}
\end{quote}

\begin{quote}
Do you go along with my view?
\end{quote}

A feminine noun.

\begin{cousinweb}
\textbf{\pa{ਰਾਏ}} (\emph{rāe}) is a loan, usually traced through Persian to Arabic \emph{raʾy}, \textquotedblleft{}opinion\textquotedblright{}.
\end{cousinweb}

\begin{grammarlens}[title={What you've built}]
One more word for asking and answering.
\end{grammarlens}

\subsection*{Your turn}
\begin{itemize}
\raggedright
  \item \emph{Say it:} \emph{rāe}
  \item \emph{Say it:} \emph{rāe}, once more
  \item \emph{Say it:} \emph{kī tusī̃ merī rāe mannde ho?} — Do you go along with my view?
  \item \emph{Recall:} say the Punjabi for a condition, then the Punjabi for a decision, then say \emph{rāe} again
  \item \emph{Write it:} copy \textbf{\pa{ਰਾਏ}} once, checking the whole word
\end{itemize}

\begin{tcolorbox}[breakable,colback=teal!4,colframe=teal!35!black,title={Before you move on}]
What does \pa{ਰਾਏ} mean? (\textquotedblleft{}An opinion; a view\textquotedblright{}.) Say it once more.
\end{tcolorbox}
\section[\texorpdfstring{\pa{ਮੁਕਾਉਣਾ} (mukāuṇā)}{ਮੁਕਾਉਣਾ (mukāuṇā)}]{\pa{ਮੁਕਾਉਣਾ} (mukāuṇā) — to settle a matter; to bring to an end}

\label{lesson:PA-C345-mukauna}



\begin{quote}
Before the new one: say the Punjabi for a decision, then the Punjabi for an opinion.
\end{quote}

\subsection*{You'll want to know: \pa{ਮੁਕਾਉਣਾ}}
\textbf{\pa{ਮੁਕਾਉਣਾ}} — \emph{mukāuṇā} — \textquotedblleft{}to settle a matter; to bring to an end\textquotedblright{}.

\begin{quote}
\pa{ਕੀ} \pa{ਅਸੀਂ} \pa{ਅੱਜ} \pa{ਗੱਲ} \pa{ਮੁਕਾਵਾਂਗੇ}?
\end{quote}

\begin{quote}
\emph{kī asīṁ ajj gall mukāvāṅge?}
\end{quote}

\begin{quote}
Shall we settle the matter today?
\end{quote}

Past (I …): \pa{ਮੈਂ} \pa{ਮਾਮਲਾ} \pa{ਮੁਕਾਇਆ}, \emph{maiṁ māmlā mukāiā}.

\begin{cousinweb}
\textbf{\pa{ਮੁਕਾਉਣਾ}} (\emph{mukāuṇā}) is the causative partner of \emph{mukkṇā}, \textquotedblleft{}to come to an end\textquotedblright{}, which goes back to a reconstructed \emph{*mukna}, \textquotedblleft{}loosed\textquotedblright{}; CDIAL lists Punjabi \emph{mukkṇā}, \textquotedblleft{}to come to an end, run short\textquotedblright{}, under it (CDIAL 10157).
\end{cousinweb}

\begin{grammarlens}[title={What you've built}]
One more word for asking and answering.
\end{grammarlens}

\subsection*{Your turn}
\begin{itemize}
\raggedright
  \item \emph{Say it:} \emph{mukāuṇā}
  \item \emph{Say it:} \emph{mukāuṇā}, once more
  \item \emph{Say it:} \emph{kī asīṁ ajj gall mukāvāṅge?} — Shall we settle the matter today?
  \item \emph{Recall:} say the Punjabi for a decision, then the Punjabi for an opinion, then say \emph{mukāuṇā} again
  \item \emph{Write it:} copy \textbf{\pa{ਮੁਕਾਉਣਾ}} once, checking the whole word
\end{itemize}

\begin{tcolorbox}[breakable,colback=teal!4,colframe=teal!35!black,title={Before you move on}]
What does \pa{ਮੁਕਾਉਣਾ} mean? (\textquotedblleft{}To settle a matter; to bring to an end\textquotedblright{}.) Say it once more.
\end{tcolorbox}
\section[\texorpdfstring{\pa{ਤਜਵੀਜ਼} (tajvīz)}{ਤਜਵੀਜ਼ (tajvīz)}]{\pa{ਤਜਵੀਜ਼} (tajvīz) — a proposal; a plan put forward}

\label{lesson:PA-C345-tajviz}



\begin{quote}
Before the new one: say the Punjabi for an opinion, then the Punjabi for to settle a matter.
\end{quote}

\subsection*{You'll want to know: \pa{ਤਜਵੀਜ਼}}
\textbf{\pa{ਤਜਵੀਜ਼}} — \emph{tajvīz} — \textquotedblleft{}a proposal; a plan put forward\textquotedblright{}.

\begin{quote}
\pa{ਉਸ} \pa{ਦੀ} \pa{ਤਜਵੀਜ਼} \pa{ਬੁਰੀ} \pa{ਨਹੀਂ} \pa{ਹੈ।}
\end{quote}

\begin{quote}
\emph{us dī tajvīz burī nahīṁ hai.}
\end{quote}

\begin{quote}
His proposal is not a bad one.
\end{quote}

A feminine noun.

\begin{cousinweb}
\textbf{\pa{ਤਜਵੀਜ਼}} (\emph{tajvīz}) is a loan, usually traced through Persian to Arabic \emph{tajwīz}, \textquotedblleft{}allowing\textquotedblright{}.
\end{cousinweb}

\begin{grammarlens}[title={What you've built}]
One more word for asking and answering.
\end{grammarlens}

\subsection*{Your turn}
\begin{itemize}
\raggedright
  \item \emph{Say it:} \emph{tajvīz}
  \item \emph{Say it:} \emph{tajvīz}, once more
  \item \emph{Say it:} \emph{us dī tajvīz burī nahīṁ hai.} — His proposal is not a bad one.
  \item \emph{Recall:} say the Punjabi for an opinion, then the Punjabi for to settle a matter, then say \emph{tajvīz} again
  \item \emph{Write it:} copy \textbf{\pa{ਤਜਵੀਜ਼}} once, checking the whole word
\end{itemize}

\begin{tcolorbox}[breakable,colback=teal!4,colframe=teal!35!black,title={Before you move on}]
What does \pa{ਤਜਵੀਜ਼} mean? (\textquotedblleft{}A proposal; a plan put forward\textquotedblright{}.) Say it once more.
\end{tcolorbox}
